% 65-41(65쪽) — 도함수 y=f'(x) 의 그래프. 구간 [a, b] 에서 x축을 아홉 번 지난다(x축 위 봉우리 4개 · 아래 골 4개).
% 원문에 있는 것만: 곡선 · a, b 로 가는 세로 점선 둘 · 라벨 a, b, O, x, y, y=f'(x). 눈금·좌표는 원문에 없으므로 넣지 않는다.
% 극대·극소의 개수를 세는 문항이므로 x축을 지나는 자리와 방향만 원문 그대로 옮긴다. 스캔 화질이 낮아 TikZ 로 다시 그림.
\begin{tikzpicture}[x=0.52cm, y=0.52cm, line width=0.5pt]
  \draw[->, >=stealth, line width=0.7pt] (-5.75,0) -- (6.35,0) node[below=1pt, font=\small] {$x$};
  \draw[->, >=stealth, line width=0.7pt] (0,-2.75) -- (0,2.35) node[left=1pt, font=\small] {$y$};
  \draw[densely dashed, line width=0.4pt] (-5.13,0) -- (-5.13,-2.33);
  \draw[densely dashed, line width=0.4pt] (5.62,0) -- (5.62,1.67);
  \draw[line width=0.6pt] plot[smooth, tension=0.6] coordinates {
    (-5.13,-2.33) (-4.90,-1.55) (-4.70,-0.60) (-4.47,0) (-4.30,0.42) (-4.13,0.45)
    (-3.95,0.30) (-3.72,0) (-3.50,-0.75) (-3.30,-1.20) (-3.13,-1.33) (-2.95,-1.15)
    (-2.80,-0.70) (-2.63,0) (-2.45,0.60) (-2.30,0.95) (-2.22,1.00) (-2.05,0.85)
    (-1.85,0.45) (-1.63,0) (-1.45,-0.90) (-1.20,-1.90) (-0.97,-2.33) (-0.75,-1.95)
    (-0.50,-1.10) (-0.22,0) (0.00,1.00) (0.20,1.65) (0.37,1.83) (0.55,1.60)
    (0.75,1.00) (1.03,0) (1.20,-0.35) (1.53,-0.42) (1.80,-0.20) (2.03,0)
    (2.20,0.55) (2.40,0.80) (2.53,0.83) (2.70,0.70) (2.88,0.35) (3.03,0)
    (3.30,-0.85) (3.60,-1.35) (4.03,-1.50) (4.40,-1.20) (4.80,-0.50) (5.20,0)
    (5.45,0.90) (5.62,1.67)};
  \node[above=1pt, font=\small] at (-5.13,0) {$a$};
  \node[below=1pt, font=\small] at (5.62,0) {$b$};
  \node[below right=-1pt and -1pt, font=\small] at (0,0) {$\pt{O}$};
  \node[anchor=west, font=\small] at (3.55,1.95) {$y=f'(x)$};
\end{tikzpicture}
